% ============================================================
\chapter{Machine Learning with TDA}
\label{ch:ml_tda}
% ============================================================

Persistence diagrams capture the ``shape'' of data, but machine learning algorithms --- SVMs, random forests, neural networks --- expect fixed-dimensional vectors. How does one bridge this gap? This is the problem of \emph{vectorization}, one of the central challenges of applied TDA. Peter Bubenik (2015) proposed \emph{persistence landscapes}, functions in a Banach space; Henry Adams and collaborators introduced \emph{persistence images}; Mathieu Carri\`ere et al.\ developed persistence kernels. More recently, differentiable persistence layers allow integrating TDA directly into deep neural networks, opening the way to end-to-end topological learning.

\begin{intuition}
Persistence diagrams live in a non-vectorial space: one cannot
directly add them or multiply them by a scalar. To use them in
a machine learning pipeline, we must \textbf{vectorize}\index{vectorization}
them, i.e., transform them into representations living in a Hilbert
space or a finite-dimensional feature space.
\end{intuition}

% ============================================================
\section{Persistence landscapes}
\label{sec:landscapes}
% ============================================================

Recall the persistence landscape introduced in
Chapter~\ref{ch:diagrams}: for a diagram $D = \{(b_i, d_i)\}$, the
$k$-th landscape $\lambda_k(t)$ is the $k$-th largest value of the
tent functions $\Lambda_i(t)$. Landscapes live in a Banach space,
which allows computing means, variances, and performing statistical
tests.
\index{persistence landscape}

\begin{minted}{python}
from gudhi.representations import Landscape
import numpy as np

# Persistence diagram
dgm = np.array([[0.0, 1.5], [0.5, 2.0], [1.0, 3.0]])

# Compute first 3 landscapes on a grid
landscape = Landscape(num_landscapes=3, resolution=100)
L = landscape.fit_transform([dgm])
print(f"Shape: {L.shape}")  # (1, 300)
\end{minted}

% ============================================================
\section{Persistence images}
\label{sec:images}
% ============================================================

\begin{definition}[Persistence image --- Adams et al.\ 2017]
\index{persistence image}
Let $D$ be a persistence diagram. The \textbf{persistence image}
is constructed in three steps:
\begin{enumerate}
  \item \textbf{Coordinate change}: transform each $(b_i, d_i)$
        into $(b_i, d_i - b_i)$ (birth, persistence).
  \item \textbf{Weighting}: apply a weight function
        $w : \R^2 \to \R_+$ (typically linear in persistence)
        to attenuate points close to the diagonal.
  \item \textbf{Gaussian smoothing}: place a Gaussian
        $\mathcal{N}((b_i, p_i), \sigma^2 I)$ on each point and
        discretize on an $N \times N$ grid.
\end{enumerate}
The resulting image is:
\[
  \rho(x, y) = \sum_{i=1}^n w(b_i, p_i) \cdot
  \frac{1}{2\pi\sigma^2}
  \exp\!\left(-\frac{(x - b_i)^2 + (y - p_i)^2}{2\sigma^2}\right).
\]
\end{definition}

\begin{theorem}[Stability of persistence images]
If $w$ is Lipschitz and bounded, the map $D \mapsto \rho_D$ is
Lipschitz with respect to the Wasserstein distance $W_1$:
\[
  \norm{\rho_{D_1} - \rho_{D_2}}_\infty \leq C \cdot W_1(D_1, D_2).
\]
\end{theorem}

% ============================================================
\section{Persistence silhouettes}
\label{sec:silhouettes}
% ============================================================

Recall the persistence silhouette of order $q$ introduced in
Chapter~\ref{ch:diagrams}: it is the weighted average (by
$(d_i - b_i)^q$) of the tent functions. For large $q$, it gives more
weight to persistent features. The silhouette is an element of
$L^p(\R)$, hence directly usable as a feature in a machine learning
pipeline.
\index{persistence silhouette}

% ============================================================
\section{Kernel methods}
\label{sec:kernels}
% ============================================================

\begin{definition}[Persistence kernel]
\index{persistence kernel}
A \textbf{persistence kernel} is a positive definite kernel
$K : \mathcal{D} \times \mathcal{D} \to \R$ on the space of
persistence diagrams.
\end{definition}

\begin{definition}[Persistence scale kernel --- Reininghaus et al.\ 2015]
\[
  k_\sigma(D_1, D_2) = \frac{1}{8\pi\sigma}
  \sum_{p \in D_1} \sum_{q \in D_2}
  \left(e^{-\frac{\norm{p - q}^2}{8\sigma}}
  - e^{-\frac{\norm{p - \bar{q}}^2}{8\sigma}}\right)
\]
where $\bar{q} = (d_q, b_q)$ is the point $q$ reflected across the
diagonal.
\end{definition}

\begin{proposition}[Positive definiteness]
The kernel $k_\sigma$ is positive definite and stable:
$\abs{k_\sigma(D_1, D_1') - k_\sigma(D_2, D_2')} \leq
C_\sigma (W_1(D_1, D_2) + W_1(D_1', D_2'))$.
\end{proposition}

\begin{definition}[Fisher kernel --- Le, Yamada 2018]
\index{Fisher kernel}
The \textbf{Fisher kernel} models each diagram as a distribution
and uses the Fisher distance between distributions:
\[
  K_{\mathrm{Fisher}}(D_1, D_2) = \exp\!\left(
  -\frac{1}{2t} d_{\mathrm{Fisher}}^2(D_1, D_2)\right).
\]
\end{definition}

% ============================================================
\section{Integration into an ML pipeline}
\label{sec:pipeline}
% ============================================================

\begin{example}[Complete classification pipeline]
\begin{enumerate}
  \item \textbf{Data}: point clouds or time series.
  \item \textbf{Filtration}: Rips, Alpha, or sublevel set.
  \item \textbf{Persistence}: compute diagrams via \texttt{gudhi}
        or \texttt{ripser}.
  \item \textbf{Vectorization}: persistence images, landscapes,
        or kernels.
  \item \textbf{Model}: SVM, Random Forest, or logistic regression.
\end{enumerate}
\end{example}

\begin{minted}{python}
from sklearn.pipeline import Pipeline
from sklearn.svm import SVC
from gudhi.representations import (
    PersistenceImage, Landscape, DiagramSelector
)

# scikit-learn pipeline with topological features
pipe = Pipeline([
    ("selector", DiagramSelector(use=True)),
    ("persistence_image", PersistenceImage(
        bandwidth=0.1, resolution=[20, 20]
    )),
    ("classifier", SVC(kernel="rbf", C=10.0)),
])

# pipe.fit(X_train_diagrams, y_train)
# y_pred = pipe.predict(X_test_diagrams)
\end{minted}

\begin{attention}[title=\textbf{Choosing a vectorization}]
\begin{itemize}
  \item \textbf{Landscapes} preserve Banach space structure and
        allow statistical tests.
  \item \textbf{Images} are better suited for CNNs since they
        produce 2D representations.
  \item \textbf{Kernels} are preferable with SVMs in low dimension.
\end{itemize}
\end{attention}

% ============================================================
\section{Scikit-TDA and software ecosystem}
\label{sec:scikit_tda}
% ============================================================

\begin{remark}[Python ecosystem for TDA]
\index{scikit-TDA}
\begin{itemize}
  \item \textbf{GUDHI}: complete C++/Python library (filtrations,
        persistence, representations).
  \item \textbf{Ripser}: fast Rips persistence computation.
  \item \textbf{scikit-tda}: meta-package integrating Ripser,
        Persim (distances), KeplerMapper.
  \item \textbf{giotto-tda}: TDA pipeline compatible with
        scikit-learn, with C++ acceleration.
\end{itemize}
\end{remark}

% ============================================================
\section{Exercises}
% ============================================================

\begin{exercise}
Compute the first two persistence landscapes of the diagram
$D = \{(0, 3), (1, 4), (2, 5)\}$ and plot them.
\end{exercise}

\begin{exercise}
Show that the persistence silhouette of order $q = 1$ is the
persistence-weighted average of the tent functions.
\end{exercise}

\begin{exercise}[Implementation]
Build a complete pipeline that classifies point clouds sampled
from a circle vs.\ a figure eight, using persistence images
and an SVM.
\end{exercise}

\begin{exercise}
Show that the Reininghaus kernel is symmetric:
$k_\sigma(D, D') = k_\sigma(D', D)$.
\end{exercise}

\begin{exercise}[Experimental comparison]
Compare classification performance (accuracy, computation time)
across landscapes, images, and persistence kernels on a dataset
of your choice.
\end{exercise}

\begin{exercise}[Persistence landscape stability]
Let $D_1$ and $D_2$ be two persistence diagrams and
$\lambda_k^{(1)}$, $\lambda_k^{(2)}$ their respective $k$-th landscapes.
\begin{enumerate}
  \item Show that for each diagram point $(b, d)$, the tent function
        $\Lambda(t) = \max(0, \min(t - b, d - t))$ is $1$-Lipschitz.
  \item Using the fact that the $\mathrm{kmax}$ operation is
        $1$-Lipschitz in $\ell^\infty$ norm on sorted vectors, show that:
        \[
          \norm{\lambda_k^{(1)} - \lambda_k^{(2)}}_{L^\infty}
          \leq W_\infty(D_1, D_2).
        \]
  \item Generalize to the $L^p$ case: show that
        $\norm{\lambda_k^{(1)} - \lambda_k^{(2)}}_{L^p}
        \leq C_p \cdot W_p(D_1, D_2)^{1/p}$
        and discuss the constant $C_p$.
\end{enumerate}
\end{exercise}

\begin{exercise}[Kernel SVM pipeline with persistence features]
The goal is to classify point clouds into three classes: circle,
torus (2D cross-section), and sphere (sampled in 3D).
\begin{enumerate}
  \item For each point cloud, compute the persistence diagram in
        homology degrees 0 and 1 via a Vietoris--Rips complex.
  \item Build the Gram matrix of the Reininghaus kernel:
        $G_{ij} = k_\sigma(D_i, D_j)$ for a parameter $\sigma$ chosen
        by cross-validation.
  \item Train a kernel SVM with precomputed kernel
        (\texttt{kernel='precomputed'}) and report accuracy under
        5-fold cross-validation.
  \item Compare with a pipeline using concatenated persistence images
        (degrees 0 and 1) and an RBF SVM. Discuss the advantages and
        disadvantages of each approach.
\end{enumerate}
\end{exercise}
